% !TEX root = ../main.tex
\part{Exércitos}

\section{Construção de Pelotão}

Vanguarda 44 combina dois sistemas de construção de força, cada um otimizado para sua escala: o Esquadrão de Infantaria usa os Slots de Dado (Parte IV, sem contagem de pontos), enquanto o Pelotão como um todo --- oficiais, número de esquadrões, suporte e veículos --- usa um sistema de pontos simples, no espírito do Pelotão Reforçado do Bolt Action.

\subsection{Estrutura Mínima}

Todo Pelotão precisa de, no mínimo: 1 Oficial (Tenente ou superior) e 2 Esquadrões de Infantaria.

\subsection{Slots de Pelotão (Sugestão de Organização)}

\begin{center}
\begin{tabular}{p{2.6cm}p{5.2cm}p{4.6cm}}
\toprule
Categoria & Exemplos & Quantidade Sugerida \\
\midrule
Comando & Tenente + Comando de Pelotão & 1 obrigatório \\
Linha de Frente & Esquadrões de Infantaria (10 dados cada) & 2 obrigatórios, +0-3 opcionais \\
Suporte & HMG, Morteiro Leve, Sniper, Médico independente & 0-2, conforme os PC disponíveis do Tenente \\
Artilharia / Blindados & Morteiro Pesado, Canhão AT, Tanque, Veículo Blindado & 0-1 a 0-2, dependendo do tamanho do jogo \\
\bottomrule
\end{tabular}
\end{center}

\subsection{Modelo de Lista de Exército (Custo em Pontos)}
\label{sec:tabela-pontos}

Para o nível de Pelotão, use custo em pontos por unidade --- não por miniatura individual dentro do esquadrão:

\begin{center}
\begin{tabular}{p{1.8cm}p{4.6cm}p{1.6cm}p{1.4cm}p{1.4cm}}
\toprule
Tipo & Nome da Unidade & Pontos & Mín. & Máx. \\
\midrule
QG & Comando de Pelotão (Oficial + Rádio/Ajudante) & 70 & 1 & 3 \\
Infantaria & Esquadrão de Fuzileiros & 65 & 5 & 10 \\
Suporte & Equipe LMG / HMG & 50 & 2 & 3 \\
Suporte & Equipe de Sniper & 60 & 2 & 2 \\
Veículo & Tanque Leve & 150 & 1 & 1 \\
\bottomrule
\end{tabular}
\end{center}

\begin{nota}
Uma partida introdutória costuma funcionar bem entre 500 e 750 pontos; partidas padrão, entre 1000 e 1250 pontos --- valores de referência, ajustáveis conforme o grupo.
\end{nota}

\section{Doutrinas Nacionais e Facções}

Cada exército lutou a Segunda Guerra de maneira diferente, focado em seu próprio equipamento e filosofia de combate. Ao montar sua lista, escolha seu Exército e, em seguida, sua Sub-Facção.

\subsection{Estados Unidos da América}

\textbf{Doutrina Nacional: Fogo Superior e Mobilidade}

Os americanos tinham o M1 Garand (fuzil semiautomático) e uma logística impecável, permitindo atirar mais rápido e se mover constantemente.

\textbf{Regra do Exército (Fogo em Movimento)}: Fuzileiros americanos NÃO sofrem a penalidade de -1 para acertar quando recebem a ordem de Avançar.

\subsubsection{Sub-Facção A: Exército Regular (G.I.s)}

Foco: Artilharia e Comunicação. Regra Especial (``Rede de Rádio''): recebem 1 Operador de Rádio gratuito para o Tenente (não conta no limite de modelos do QG), e podem rolar novamente o dado de Ajuste de Mira de Morteiros que tirem resultado 1.

\subsubsection{Sub-Facção B: US Airborne (Paraquedistas)}

Foco: Elite, lutar cercado e resiliência. Regra Especial (``Band of Brothers''): Tropas Airborne ignoram os 2 primeiros níveis de Stress que receberiam a cada turno (1 do bônus padrão de Veterano + 1 da doutrina).

\subsection{Alemanha (Wehrmacht)}

\textbf{Doutrina Nacional: A Máquina de Serrar}

A doutrina de infantaria alemã não girava em torno do fuzileiro, mas da Metralhadora (MG34/MG42). Todo o esquadrão existia para manter a MG atirando.

\textbf{Regra do Exército}: todas as LMGs e HMGs alemãs recebem +1 Dado de Ataque gratuito. Porém, se o Municiador da arma morrer, a penalidade de Operação Solo é de -2 no Teste de Acerto (em vez do -1 padrão).

\subsubsection{Sub-Facção A: Heer (Exército Regular)}

Foco: Disciplina de oficiais e flexibilidade tática. Regra Especial (``Auftragstaktik''): se o Sargento de um esquadrão for morto, a unidade não sobe nenhum nível no dado de Stress pela perda do líder --- um fuzileiro é imediatamente ``promovido'', garantindo que o esquadrão continue gerando seu 1 PC no turno seguinte.

\subsubsection{Sub-Facção B: Fallschirmjäger / Waffen-SS}

Foco: Fanatismo e combate de curta distância. Regra Especial (``Até o Último Homem''): ao atingir o Nível 6 de Stress (que normalmente causaria fuga), role 1D6: com um resultado de 4+, a unidade permanece em campo --- sofrendo todas as penalidades do Nível 6 --- em vez de ser removida, até reduzir ao menos 1 nível de Stress.
